% Préambule commun aux deux documents
\usepackage[T1]{fontenc}
\usepackage[french]{babel}
\usepackage{lmodern}
\usepackage{graphicx,xcolor,booktabs,array,tabularx,longtable}
\usepackage{tikz}
\usetikzlibrary{positioning,arrows.meta,shapes.geometric,calc,fit,backgrounds,decorations.pathreplacing}
\usepackage{fancyhdr,enumitem,microtype}
\usepackage[colorlinks=true,linkcolor=insared,urlcolor=insared,citecolor=insared]{hyperref}
\usepackage{caption}
\usepackage{listings}

\definecolor{insared}{RGB}{200,16,46}
\definecolor{grisfond}{RGB}{247,248,250}
\definecolor{grisligne}{RGB}{216,220,226}
\definecolor{gristexte}{RGB}{91,97,107}
\definecolor{vertok}{RGB}{26,127,75}
\definecolor{orangepartiel}{RGB}{178,106,5}
\definecolor{rougeko}{RGB}{198,40,40}

\captionsetup{font=small,labelfont={bf,color=insared}}
\setlist[itemize]{leftmargin=1.2em,itemsep=1pt,topsep=2pt,parsep=0pt}
\setlist[enumerate]{leftmargin=1.4em,itemsep=1pt,topsep=2pt,parsep=0pt}

\lstset{
  basicstyle=\ttfamily\scriptsize,
  backgroundcolor=\color{grisfond},
  frame=single, rulecolor=\color{grisligne},
  breaklines=true, showstringspaces=false,
  keywordstyle=\color{insared}\bfseries,
  commentstyle=\color{gristexte}\itshape,
  numberstyle=\tiny\color{gristexte},
  xleftmargin=4pt, framexleftmargin=4pt, aboveskip=6pt, belowskip=6pt,
  literate={é}{{\'e}}1 {è}{{\`e}}1 {à}{{\`a}}1 {ç}{{\c c}}1 {ê}{{\^e}}1 {û}{{\^u}}1 {î}{{\^i}}1 {ô}{{\^o}}1 {É}{{\'E}}1 {—}{{---}}1 {’}{{'}}1 {«}{{<<}}1 {»}{{>>}}1 {·}{{\textperiodcentered}}1
}

% Marqueurs de conformité
\newcommand{\ok}{\textcolor{vertok}{\textbf{Conforme}}}
\newcommand{\partiel}{\textcolor{orangepartiel}{\textbf{Partiel}}}
\newcommand{\absent}{\textcolor{rougeko}{\textbf{Absent}}}
\newcommand{\plus}{\textcolor{insared}{\textbf{Au-delà}}}

% Les paragraphes techniques enchaînent de longues suites en \texttt{} que TeX
% ne peut pas couper : sans marge de man\oe{}uvre, elles débordent dans la marge.
% \emergencystretch autorise un dernier passage plus souple sur ces lignes-là.
\emergencystretch=2em
\hbadness=4000        % ne pas signaler les lignes simplement peu denses
